* PLOT: Scheme of the benchmark

Our benchmark agregates datasets with each dataset preprocessed with the same pipeline, with the exception of HCP which has its own preprocessing pipeline. We perform the feature extraction/data preparation and apply rigorously the analysis and cross validation that we are going to explain.
\subsection{Datasets}
\begin{itemize}
    \item Number of subjects
    \item Age range / demographics
    \item Scanner details
    \item Acquisition protocol or summary table with acquisition parameters
    \item Prediction targets (classification/regression)
\end{itemize}
* PLOTS: Demographic distribution (age histograms) | Table with demographics summary (gender classes, age range...)

\subsection{Preprocessing and Feature Extraction}
\subsubsection{Preprocessing (Himanshu)}
\begin{itemize}
    \item Denoising
    \item Motion and eddy correction
    \item Spatial normalization
    \item Masking strategy
\end{itemize}

\subsubsection{Feature extraction}
\begin{itemize}
    \item White and Gray matter computation (Maybe merge in Matter representation and include a table with differences between both representations(?))
    \item Input: Microstructure information
\end{itemize}
    \paragraph{Gray Matter Representation}
    \begin{itemize}
        \item Voxel-wise features
        \item High-dimensional representation (~64k features)
        \item Mask-based extraction
    \end{itemize}
    Explain that voxel-wise features preserve spatial resolution but increase dimensionality and risk of overfitting.

    \paragraph{White Matter Representation}
    \begin{itemize}
        \item Tract-based aggregation
        \item Mean values per tract (~48 features)
        \item Reduced dimensionality
        \item Anatomical interpretability
    \end{itemize}
    Justify tract-averaged representation: Noise reduction, improved stability, lower dimensional feature space, anatomically meaningful structure

\subsection{Prediction}
\begin{itemize}
    \item Feature normalization
        \begin{itemize}
            \item Binary targets (encoded 0/1), Continuous targets (standardized during training)
            \item Standardization fitted on training set only
            \item Applied to validation/test sets 
            \item Inverse transform for regression outputs (for visualization)
        \end{itemize}
    \item Task definitions and labels and reasons for task selection
        \begin{itemize}
            \item Classification: Binary diagnosis prediction (e.g., ASD vs. TD)
            \item Regression: Age prediction (brain age estimation)
            \item Task selection rationale: Clinical relevance, data availability, and established benchmarks in the literature
        \end{itemize}
    \item Evaluation Metrics
        \begin{itemize}
            \item Classification: Accuracy, AUC, F1-score
            \item Regression: R², MAE
            \item Mean ± std across folds
            \item Statistical comparison tests (e.g., paired t-tests, Wilcoxon signed-rank tests) to assess significance of performance differences between models
        \end{itemize}
    \item Models Used
        \begin{itemize}
            \item Dummy baseline (e.g., majority class for classification, mean predictor for regression)
            \item Classical ML pipeline: Linear models (e.g., logistic regression, linear regression), Random Forests, Support Vector Machines (SVM) + their PCA versions
            \item Deep Learning models (e.g., dinov2, medicalnet)
            \item Unified hyperparameter search strategy (e.g., grid search or random search with cross-validation on the training set)
        \end{itemize}
\end{itemize}

\subsection{Validation, Evaluation and Visualizations}
\begin{itemize}
    \item Data Splitting Strategy (Fixed train/test splits)
    \item Cross-validation within training set. Stratification
    \item No information leakage
    \item Explanation of plots and statistics
\end{itemize}